
\chapter*{\MakeUppercase{Kata Pengantar}}
\addcontentsline{toc}{chapter}{\MakeUppercase{Kata Pengantar}}

Kata pengantar merupakan bagian awal tugas akhir yang berfungsi memberikan gambaran singkat mengenai penyusunan tugas akhir. Pada bagian ini, penulis menjelaskan latar belakang penyusunan dokumen, tujuan penulisan, serta harapan terhadap manfaat yang dapat diberikan oleh hasil penelitian yang dilakukan.

Kata pengantar ditulis menggunakan bahasa yang formal, ringkas, dan mudah dipahami. Bagian ini tidak digunakan untuk memaparkan hasil penelitian, analisis, maupun kesimpulan karena pembahasan tersebut telah disajikan pada bab-bab utama tugas akhir.

Secara umum, kata pengantar dapat memuat hal-hal berikut:

\begin{enumerate}[topsep=0pt,itemsep=0pt,parsep=0pt,leftmargin=*]
\item Ungkapan puji syukur kepada Tuhan Yang Maha Esa atas terselesaikannya tugas akhir.
\item Penjelasan singkat mengenai tujuan penyusunan tugas akhir.
\item Gambaran umum mengenai topik atau ruang lingkup penelitian yang dilakukan.
\item Harapan penulis terhadap manfaat hasil penelitian bagi pembaca, pengembangan ilmu pengetahuan, maupun penerapan di bidang terkait.
\item Pernyataan mengenai keterbatasan penulisan serta keterbukaan terhadap kritik dan saran yang membangun.
\end{enumerate}

Hal-hal yang perlu diperhatikan dalam penulisan kata pengantar adalah sebagai berikut:

\begin{itemize}[topsep=0pt,itemsep=0pt,parsep=0pt,leftmargin=*]
\item Ditulis secara ringkas dan tidak berlebihan.
\item Menggunakan bahasa yang formal dan santun.
\item Tidak memuat ucapan terima kasih kepada pihak tertentu apabila telah disediakan bagian \textit{Ucapan Terima Kasih} secara terpisah.
\item Tidak memuat hasil penelitian, analisis, maupun kesimpulan penelitian.
\item Panjang kata pengantar umumnya tidak lebih dari satu halaman.
\end{itemize}
